%----------------------------------------------------------------------------
\chapter{Weekly consultations}\label{app:consult}
%----------------------------------------------------------------------------
\paragraph{Find a recurring slot.} Agree on a weekly slot with us and create a recurring MS Teams event.

\paragraph{Why every week.} A week is short enough that a wrong turn costs days, not a month, and long enough to have something to show. Come even when little has moved: a blocked week is exactly the one we can help with, and the one students tend to skip.

\paragraph{Why one-on-one.} A group meeting averages over everyone's problem, and yours is specific: your code, your draft, your plot. Alone with you we read the actual artifact and ask back until the misunderstanding is found, which takes minutes one-on-one and rarely happens in a group. It also lets us be frank about scope and effort, which is not a conversation for an audience.

\paragraph{Why in person.} In a room we share a whiteboard, which is where notation, architecture and the leading figure actually get drawn. We also see where you hesitate, which a call hides, and the hard conversations, when a topic has to be cut or re-aimed, go better face to face.

\paragraph{Consultation slides are optional.} A few slides on what you did this week, or slides that make your work easier for us to understand, save the first ten minutes of every consultation; we welcome them, LLM-generated included. A useful outline:
\begin{enumerate}[noitemsep]
  \item what you did since the last consultation, with evidence: a figure, a result, a page of text;
  \item what you read: the papers you added to Zotero this week, one sentence each on why they matter;
  \item what is blocked or where you are unsure;
  \item what you plan for next week;
  \item the questions you need us to decide.
\end{enumerate}

\paragraph{Close with a summary.} End by saying back in your own words what we agreed: the decisions, your next steps and what we owe you. Then save the same as a short written summary in your notes repository, as \texttt{notes/consultations/YYYY-MM-DD*.md}. If your summary differs from what we meant, we catch it now rather than a week of work later. The summaries are also where anchor approvals and your decisions on reviewer objections are recorded, so they are the paper's memory as much as ours.

\paragraph{The mirror principle.} We put in what you put in, and we are open about it. A student who works hard on their topic gets a lot of our time, our contacts and our help beyond the paper; a student who has put the topic on the back burner gets less of it. This is not about struggling: a student who struggles and keeps showing up gets more help, not less. It means we are mindful of where our time goes, and that you decide how much of it you get. On a separate matter, we value consistency over bursts. A student who is diligent every week and not very productive gets further than one who disappears for a month and then does brilliant work for a weekend, because research compounds: each week's result is what the next week builds on, and a weekly rhythm keeps the topic alive in your head and in ours.

\paragraph{Recordings and Teams calls.} Hold the consultation as a Teams call even when we sit in one room, and record it with transcription. You can then rewatch the details, feed the transcript to your AI tools and turn it into the notes in your repository. We recommend \href{https://github.com/openai/whisper}{Whisper} large-v3-turbo over Teams' built-in transcription: it is more accurate and handles multilingual speech better, and we usually speak a mix of Hungarian and English. Ask us for our transcription skill.
